
\documentclass{article}
\usepackage{colortbl}
\usepackage{makecell}
\usepackage{multirow}
\usepackage{supertabular}

\begin{document}

\newcounter{utterance}

\centering \large Interaction Transcript for game `cladder', experiment `full\_v1.5\_default', episode 9771 with Qwen3.5{-}9B.
\vspace{24pt}

{ \footnotesize  \setcounter{utterance}{1}
\setlength{\tabcolsep}{0pt}
\begin{supertabular}{c@{$\;$}|p{.15\linewidth}@{}p{.15\linewidth}p{.15\linewidth}p{.15\linewidth}p{.15\linewidth}p{.15\linewidth}}
    \# & \multicolumn{2}{c}{Player} && \multicolumn{2}{c}{Game Master} \\
    \hline

    \theutterance \stepcounter{utterance}  
    & & & \multicolumn{4}{p{0.6\linewidth}}{
        \cellcolor[rgb]{0.9,0.9,0.9}{
            \makecell[{{p{\linewidth}}}]{
                \texttt{\tiny{[P1$\langle$GM]}}
                \texttt{You are an expert in causal inference. The following question is not a typical commonsense query, but rather a meticulously designed question created by a professor specializing in causal inference, intended to assess the students' mastery of the course content.} \\
\texttt{Imagine a self{-}contained, hypothetical world with only the following conditions, and without any unmentioned factors or causal relationships: Jyka has a direct effect on yupt and kwox. Yupt has a direct effect on kwox. We know that jyka causes yupt. jyka or yupt causes kwox. Would an individual is kwox if not jyka instead of jyka?} \\
\texttt{Start your answer with 'Yes' or 'No', followed by additional reasoning or evidence to support your explanation.} \\
            }
        }
    }
     \\ \\

    \theutterance \stepcounter{utterance}  
    & \multicolumn{4}{p{0.6\linewidth}}{
        \cellcolor[rgb]{0.561,0.737,0.561}{
            \makecell[{{p{\linewidth}}}]{
                \texttt{\tiny{[P1$\rangle$GM]}}
                \texttt{No} \\
\\ 
\texttt{**Reasoning:**} \\
\\ 
\texttt{To determine the answer, we must analyze the causal structure and the counterfactual query within the constraints of the hypothetical world.} \\
\\ 
\texttt{1.  **Causal Structure Identification**:} \\
\texttt{    *   **Jyka \$\textbackslash{}rightarrow\$ Yupt**: The prompt explicitly states "jyka causes yupt."} \\
\texttt{    *   **Jyka \$\textbackslash{}rightarrow\$ Kwox**: The prompt states "Jyka has a direct effect on... kwox."} \\
\texttt{    *   **Yupt \$\textbackslash{}rightarrow\$ Kwox**: The prompt states "Yupt has a direct effect on kwox."} \\
\texttt{    *   **Redundancy Check**: The phrase "jyka or yupt causes kwox" is consistent with the established direct paths (\$Jyka \textbackslash{}rightarrow Kwox\$ and \$Yupt \textbackslash{}rightarrow Kwox\$). It does not imply an exclusive OR (XOR) logic that would negate the simultaneous existence of both paths; rather, it describes the set of sufficient causes for Kwox.} \\
\\ 
\texttt{2.  **Counterfactual Analysis**:} \\
\texttt{    *   **Factual World (\$W\_\{actual\}\$)**: Jyka occurs. Consequently, Yupt occurs (due to \$Jyka \textbackslash{}rightarrow Yupt\$). Both Jyka and Yupt occur, leading to Kwox occurring via two distinct direct pathways.} \\
\texttt{    *   **Counterfactual World (\$W\_\{counterfactual\}\$)**: We intervene to set Jyka to "not occur" (\$do(Jyka=0)\$).} \\
\texttt{        *   Since Jyka is the direct cause of Yupt, setting Jyka to 0 implies Yupt will also be 0 (assuming no other unmentioned causes for Yupt, which fits the "self{-}contained" constraint).} \\
\texttt{        *   Since Jyka is a direct cause of Kwox, removing Jyka removes one pathway to Kwox.} \\
\texttt{        *   However, Yupt is *also* a direct cause of Kwox. In the factual world, Yupt was active because Jyka was active. In the counterfactual world, if we simply remove Jyka, Yupt becomes inactive, and thus the pathway \$Yupt \textbackslash{}rightarrow Kwox\$ is also severed.} \\
\\ 
\texttt{    *   **Crucial Distinction (Total vs. Partial Effect)**:} \\
\texttt{        The question asks: "Would an individual {[}be{]} kwox if not jyka instead of jyka?"} \\
\texttt{        This is asking for the outcome of the variable \$Kwox\$ under the intervention \$do(Jyka=0)\$.} \\
\\ 
\texttt{        *   Path 1 (\$Jyka \textbackslash{}rightarrow Kwox\$): Broken because Jyka is removed.} \\
\texttt{        *   Path 2 (\$Jyka \textbackslash{}rightarrow Yupt \textbackslash{}rightarrow Kwox\$): Broken because the root cause (Jyka) is removed, preventing Yupt from occurring.} \\
\\ 
\texttt{        Since J} \\
            }
        }
    }
    & & \\ \\

    \theutterance \stepcounter{utterance}  
    & & & \multicolumn{2}{p{0.3\linewidth}}{
        \cellcolor[rgb]{0.9,0.9,0.9}{
            \makecell[{{p{\linewidth}}}]{
                \texttt{\tiny{[GM$|$GM]}}
                \texttt{No} \\
            }
        }
    }
    & & \\ \\

    \theutterance \stepcounter{utterance}  
    & & & \multicolumn{2}{p{0.3\linewidth}}{
        \cellcolor[rgb]{0.9,0.9,0.9}{
            \makecell[{{p{\linewidth}}}]{
                \texttt{\tiny{[GM$|$GM]}}
                \texttt{no} \\
            }
        }
    }
    & & \\ \\

    \theutterance \stepcounter{utterance}  
    & & & \multicolumn{2}{p{0.3\linewidth}}{
        \cellcolor[rgb]{0.9,0.9,0.9}{
            \makecell[{{p{\linewidth}}}]{
                \texttt{\tiny{[GM$|$GM]}}
                \texttt{game\_result = WIN} \\
            }
        }
    }
    & & \\ \\

\end{supertabular}
}

\end{document}
